\section{References and Back Matter}

The reference list is rendered below through the \texttt{ieeetr} bibliography
style from \texttt{references.bib}, and includes the five daraxonrasib program
entries; the three main documents on which this protocol builds; the directly
relevant author works; the clinical references; and the FDA guidance, the
applicable 21 CFR parts (11, 50, and 812), and the IEC, ASME, NASA, and IEEE
consensus standards together with the NIH protocol template. The three main
documents are cited here with a brief explanation of the role each plays.

\begin{itemize}[leftmargin=1.6em,itemsep=3pt,topsep=2pt]
\item The 2030 PDAC sixty-second robotic Whipple and daraxonrasib simulation
\cite{kawchak_2026_20196639} is the simulation source from which the clinical
figures, the exemplar PAT-PDAC-0001, and the operative and telemetry behavior in
this protocol derive; it supplies the in silico evidence that the early-feasibility
device arm is built upon.
\item The Adaption of 21 CFR Part 312 for end-to-end Physical AI oncology trial
unification \cite{kawchak2026adapt312} is the regulatory source for the Physical AI
overlay, including Subpart J, the $\S$312.60(f) Physical AI consent and opt-out,
the Phase 0 simulation-validation gate, the USL readiness rating, and the
hash-chained audit-trail requirements adopted throughout.
\item The H. R. 9510 bill, Verification Before Generation in Physical AI Oncology
Trials Act of 2026 \cite{kawchak2026hr9510}, is the legislative and financial
source that ties the funding and oversight of Physical AI oncology investigations
to the VVUQ evidentiary standard operationalized by the ten-gate assurance suite
(\S11.1).
\end{itemize}

\bmhead{Acknowledgments}
This protocol was assembled with Claude Code (Opus 4.8, 1M context). The author
thanks the open-source robotics, simulation, and clinical-standards communities.

\bmhead{Author and ORCID}
Kevin Kawchak\,\orcidicon{0009-0007-5457-8667}, Chief Executive Officer,
ChemicalQDevice. ORCID iD: \href{https://orcid.org/0009-0007-5457-8667}{https://orcid.org/0009-0007-5457-8667}.
Correspondence: \href{mailto:kevink@chemicalqdevice.com}{kevink@chemicalqdevice.com}.

\bmhead{Ethical Disclosures}
This is an independent research draft. No real patients were enrolled, no protected
health information was used, and no IRB approval was sought or obtained. The
exemplar PAT-PDAC-0001 is synthetic and does not represent an enrolled individual.
All clinical figures derive from the author's simulation sources and are
illustrative unless tied to a cited reference. This document is not an active IND
or IDE, is not an approved protocol, and is not medical or regulatory advice; it is
not endorsed by the FDA, NIH, HHS, an IRB, ICH, or any sponsor.

\bmhead{Data Availability}
The protocol, the analysis code, and the derived data sets are available through a
public GitHub repository and archived on Zenodo under the Creative Commons
Attribution 4.0 International (CC BY 4.0) license, at DOI
\href{https://doi.org/10.5281/zenodo.xxxxxxxx}{10.5281/zenodo.xxxxxxxx} (the
placeholder is filled at deposit). All deposited data are simulation-derived,
generated under a deterministic seed, and contain no personally identifiable
information.

\bmhead{Cite This Article}
Kawchak K. A Phase 1, First-in-Human, Combined IND/IDE Clinical Trial Protocol of
On-Premises Large Language Model Directed Robotic Pancreaticoduodenectomy (Whipple)
with Perioperative Daraxonrasib (RMC-6236) in KRAS-Mutated Pancreatic Ductal
Adenocarcinoma. Zenodo. 2026; 10.5281/zenodo.xxxxxxxx.

% Bibliography (ieeetr); \nocite{*} emits the full assembled reference list.
\phantomsection\addcontentsline{toc}{section}{References}
\begingroup\raggedright\nocite{*}\bibliographystyle{ieeetr}\bibliography{references}\endgroup
